% Einheit 1 von 3 – Kreisumfang (bank/kreis/e1.jsonl, zusammenbau v0.1)
%% TODO Zweigzeile Teil 1: was hier gelernt wird (Satz in Schülersprache aus dem Einheitstitel) – Einheit 1
\einheitenkopf[e1]{Einheit 1 von 3 · Kreisumfang}
\zweigzeile{ab Kl. 7, je nach Buch bis Kl. 8 · P10}
\setcounter{aufgabe}{8}

%% TODO Titel: Ich-kann-Satz fehlt in der Bank (Platzhalter: Kettenname) – Nr. 9
%% TODO Anweisung: ein Satz über der ersten Teilaufgabe fehlt in der Bank – Nr. 9
\begin{aufgabe}{Rand oder Fläche?}
\begin{teile}
\teil Ein runder Tisch bekommt rundherum eine Kante aus Holzleiste. Rand oder Fläche? Kreuze an.\\ \kreuz{Umfang}\\ \kreuz{Fläche}
\end{teile}
\end{aufgabe}

%% TODO Titel: Ich-kann-Satz fehlt in der Bank (Platzhalter: Kettenname) – Nr. 10
%% TODO Anweisung: ein Satz über der ersten Teilaufgabe fehlt in der Bank – Nr. 10
\begin{aufgabe}{Umfang}
%% TODO \rechenplatz{n}: Schrittzahl des Grundfalls fehlt in der Bank (nur bei fester Schreibform, 2.3 b)
\begin{teile}
\teil Ist die beschriftete Strecke ein Radius oder ein Durchmesser? Gib auch die andere Länge an.

\begin{kreis}[2] \mittelpunkt{M} \durchmesser{20}{$9$ cm} \end{kreis}
Die Strecke ist ein \leerfeld . Die andere Länge: \leerfeld
\teil Kreis mit $d = 6$ cm – wie lang ist der Umfang? Rechne mit der $\pi$-Taste und runde auf eine Stelle. u ≈ \leerfeld[cm]
\teil Kreis mit $r = 4$ cm – wie lang ist der Umfang? Rechne mit der $\pi$-Taste und runde auf eine Stelle. u ≈ \leerfeld[cm]
\teil Kreis mit $d = 4{,}6$ cm – wie lang ist der Umfang? Rechne mit der $\pi$-Taste und runde auf zwei Stellen. u ≈ \leerfeld[cm]
\teil Ein Kreis hat den Umfang $u = 47$ cm. Wie groß ist sein Durchmesser? Rechne mit der $\pi$-Taste und runde auf eine Stelle. d ≈ \leerfeld[cm]
\teil Ein Kreis hat den Umfang $u = 75$ cm. Wie groß ist sein Radius? Rechne mit der $\pi$-Taste und runde auf eine Stelle. r ≈ \leerfeld[cm]
\teil Ein Weg führt im Halbkreis um einen Teich, der Durchmesser des Halbkreises ist $d = 12$ m. Wie lang ist der Weg? Rechne mit der $\pi$-Taste und runde auf eine Stelle. ≈ \leerfeld[m]
\teil Ein Fahrradrad hat den Durchmesser $66$ cm. Wie weit fährt das Rad bei $250$ Umdrehungen? Gib das Ergebnis in Meter an, auf eine Stelle gerundet. ≈ \leerfeld[m]
\teil Ein Stifteköcher ohne Deckel hat die Form eines Zylinders mit dem Radius $3{,}6$ cm und der Höhe $11$ cm. Ein rechteckiges Papier soll genau einmal um die Seitenfläche reichen. Berechne die Länge und die Breite des Rechtecks; runde auf eine Stelle. Länge ≈ \leerfeld[cm] , Breite = \leerfeld[cm]
\end{teile}
\end{aufgabe}

%% TODO Titel: Ich-kann-Satz fehlt in der Bank (Platzhalter: Kettenname) – Nr. 11
%% TODO Anweisung: ein Satz über der ersten Teilaufgabe fehlt in der Bank – Nr. 11
\begin{aufgabe}{Radius, Durchmesser, Mittelpunkt in einer Figur benennen und einzeichnen}
\begin{teile}
\teil Zeichne in den Kreis einen Durchmesser ein, der durch den Punkt $P$ geht, und beschrifte ihn mit $d$.

\begin{kreis}[2] \mittelpunkt{M} \kreispunkt{30}{P} \end{kreis}
\end{teile}
\end{aufgabe}

%% TODO Titel: Ich-kann-Satz fehlt in der Bank (Platzhalter: Kettenname) – Nr. 12
%% TODO Anweisung: ein Satz über der ersten Teilaufgabe fehlt in der Bank – Nr. 12
\begin{aufgabe}{Kreis mit gegebenem Radius oder Durchmesser zeichnen}
\begin{teile}
\teil Zeichne einen Kreis mit dem Radius $3$ cm und markiere seinen Mittelpunkt mit $M$.

\rechenplatz[halb]{6}
\end{teile}
\end{aufgabe}

%% TODO Titel: Ich-kann-Satz fehlt in der Bank (Platzhalter: Kettenname) – Nr. 13
%% TODO Anweisung: ein Satz über der ersten Teilaufgabe fehlt in der Bank – Nr. 13
\begin{aufgabe}{Umfang messen und u : d bilden (Entdeckung von π)}
\begin{teile}
\teil An einem Becher wurde gemessen: Umfang $u = 22$ cm, Durchmesser $d = 7$ cm. Berechne $u : d$ und runde auf zwei Stellen. u : d ≈ \leerfeld
\end{teile}
\end{aufgabe}

%% TODO Titel: Ich-kann-Satz fehlt in der Bank (Platzhalter: Kettenname) – Nr. 14
%% TODO Anweisung: ein Satz über der ersten Teilaufgabe fehlt in der Bank – Nr. 14
\begin{aufgabe}{Umfang – Fehler finden, Begründen, Anwendung, Darstellungswechsel}
\begin{teile}
\teil Ein Kreis hat den Durchmesser $d = 18$ cm. Ben rechnet den Umfang: \rechnung{u &= 2 \cdot \pi \cdot 18 \\ u &\approx 113{,}1 \text{ cm}} Wo ist der Fehler?
\teil Begründe, warum man bei jedem Kreis für $u : d$ ungefähr dasselbe Ergebnis erhält.
\teil Eine Försterin misst den Umfang einer Eiche mit dem Maßband: $2{,}83$ m. Ab $75$ cm Durchmesser steht ein Baum hier unter Schutz. Steht die Eiche unter Schutz?
\teil Ergänze die Tabelle. Runde auf eine Stelle.

\wertetabelle{$d$ in cm}{$u$ in cm}{1,2,3,6}
\end{teile}
\end{aufgabe}
